\documentclass[12pt,a4paper]{article}
\usepackage{amsmath,amssymb,amsthm}
\usepackage{hyperref}
\usepackage{geometry}
\geometry{margin=2.5cm}
\usepackage{enumitem}
\usepackage{booktabs}

\newtheorem{theorem}{Theorem}[section]
\newtheorem{lemma}[theorem]{Lemma}
\newtheorem{corollary}[theorem]{Corollary}
\newtheorem{proposition}[theorem]{Proposition}
\theoremstyle{definition}
\newtheorem{definition}[theorem]{Definition}
\newtheorem{axiom_rs}[theorem]{RS Axiom}
\theoremstyle{remark}
\newtheorem{remark}[theorem]{Remark}

\title{Renormalization and Running Couplings\\
from the Recognition Science $\varphi$-Ladder}

\author{Jonathan Washburn\\
Recognition Science Research Institute\\
Austin, Texas\\
\texttt{washburn.jonathan@gmail.com}}

\date{March 2026}

\begin{document}
\maketitle

\begin{abstract}
The renormalization group (RG) describes how coupling constants run with
energy scale $\mu$: $\beta(g) = \mu \partial g/\partial\mu$.  Key
predictions include asymptotic freedom (QCD: $\alpha_s \to 0$ as $\mu
\to \infty$), the Landau pole in QED ($\alpha \to \infty$ at
$\mu \approx 10^{286}$ GeV), and gauge coupling unification near
$\mu \approx 10^{16}$ GeV.  We derive the RG structure from the
Recognition Science (RS) framework.  In RS, the coupling constants are
not fundamental — they are effective parameters emerging from the
$\varphi$-ladder at the anchor scale $\mu^\star = 182.201$ GeV (Lean:
\texttt{Physics.AnchorPolicy}).  The $\beta$-function for any coupling
$g$ is determined by the $\varphi$-ladder self-similarity: scale
transformation $\mu \to \mu e^t$ on the ledger corresponds to a rung
shift $r \to r + t/\ln\varphi$, giving $\beta_g = -g \cdot \ln\varphi
\cdot \partial\ln g/\partial r$.  Asymptotic freedom in QCD follows
from the negative $\beta$-function of the color coupling, itself derived
from the $SU(3)$ structure forced by the three-quark composition of the
proton on the $Q_3$ hypercube (Lean: \texttt{QFT.Confinement}).  The
RS anchor scale $\mu^\star$ is a stationarity point of the RG flow
(Lean: \texttt{Physics.AnchorPolicy.stationary\_at\_anchor}), not a
fit parameter.
\end{abstract}

%%============================================================
\section{Introduction}
\label{sec:intro}
%%============================================================

The renormalization group (Wilson 1971, 't Hooft 1973) describes the
scale dependence of coupling constants in quantum field theory.  For
a coupling $g(\mu)$, the $\beta$-function is
\begin{equation}
  \beta(g) = \mu \frac{\partial g}{\partial\mu}.
  \label{eq:beta}
\end{equation}

Key phenomena governed by running couplings:
\begin{itemize}
  \item \textbf{QCD asymptotic freedom}~\cite{Gross1973,Politzer1973}:
        $\beta_\mathrm{QCD} < 0$, so $\alpha_s(\mu) \to 0$ as $\mu
        \to \infty$.  This explains the success of perturbative QCD
        at high energies and quark confinement at low energies.
  \item \textbf{QED Landau pole}: $\beta_\mathrm{QED} > 0$,
        so $\alpha(\mu) \to \infty$ at the Landau pole
        $\mu_\mathrm{Landau} \sim 10^{286}$ GeV.
  \item \textbf{Gauge unification}: the three SM couplings
        $\alpha_1, \alpha_2, \alpha_3$ converge near $\mu \sim 10^{16}$
        GeV (approximately) in the MSSM.
\end{itemize}

%%============================================================
\section{Recognition Science Framework}
\label{sec:rs}
%%============================================================

\subsection{The Anchor Scale}

The RS framework fixes the anchor scale for all mass and coupling
predictions:

\begin{definition}[RS anchor scale, Lean-proved]
\label{def:anchor}
The RS anchor scale is
\begin{equation}
  \mu^\star = 182.201\,\mathrm{GeV},
  \label{eq:mustar}
\end{equation}
defined by the stationarity condition $\gamma_f(\mu^\star) = 0$
(zero anomalous dimension at the anchor).  This is proved to be a
BLM/PMS stationary point, not a fit parameter.

\texttt{Lean: Physics.AnchorPolicy.stationary\_at\_anchor}
\end{definition}

\begin{theorem}[Stationarity at anchor, Lean-proved]\label{thm:stationary}
\begin{equation}
  \frac{\partial \gamma_f(\mu^\star)}{\partial\ln\mu} = 0.
\end{equation}
\texttt{Lean: Physics.AnchorPolicy.stationary\_at\_anchor}
\end{theorem}

\begin{theorem}[Stability at anchor, Lean-proved]\label{thm:stability}
The second derivative at the anchor is bounded:
\begin{equation}
  \left|\frac{\partial^2 \gamma_f}{\partial(\ln\mu)^2}\bigg|_{\mu=\mu^\star}\right|
  \leq B
\end{equation}
for a finite bound $B$.
\texttt{Lean: Physics.AnchorPolicy.stability\_bound\_at\_anchor}
\end{theorem}

\subsection{The $\varphi$-Ladder and Scale Transformations}

In RS, a scale transformation $\mu \to \mu' = \mu e^t$ corresponds to
a rung shift on the $\varphi$-ladder:
\begin{equation}
  r \to r' = r + \frac{t}{\ln\varphi},
  \label{eq:rung_shift}
\end{equation}
since the $\varphi$-ladder is $m_r = m_0 \varphi^r$ and $\ln m_r = \ln
m_0 + r \ln\varphi$.  The RG flow is therefore translation on the
$\varphi$-ladder.

\begin{theorem}[RS $\beta$-function structure]\label{thm:beta}
For a coupling $g$ with $\varphi$-ladder dependence $g(r)$, the
RS $\beta$-function is
\begin{equation}
  \beta_g = \frac{dg}{dt} = \frac{dg}{dr}\cdot\frac{dr}{dt}
  = \frac{1}{\ln\varphi}\frac{dg}{dr}.
  \label{eq:beta_rs}
\end{equation}
The $\beta$-function sign is determined by the derivative of $g$ on the ladder.
\end{theorem}

%%============================================================
\section{Running Couplings in RS}
\label{sec:running}
%%============================================================

\subsection{QCD Strong Coupling}

The QCD $\beta$-function at one loop is
\begin{equation}
  \beta(\alpha_s) = -\frac{\alpha_s^2}{2\pi}\left(11 - \frac{2n_f}{3}\right)
  \equiv -\frac{b_0 \alpha_s^2}{2\pi},
  \label{eq:beta_qcd}
\end{equation}
where $n_f$ is the number of active quark flavors and $b_0 = 11 - 2n_f/3 > 0$
for $n_f \leq 16$.  In the SM with $n_f = 6$: $b_0 = 11 - 4 = 7$.

\begin{theorem}[Asymptotic freedom from RS]\label{thm:asym_free}
The negative $\beta$-function~\eqref{eq:beta_qcd} follows from the
RS structure of the $SU(3)$ color group:
\begin{enumerate}
  \item The factor $11$ = $11C_A/3$ comes from the gluon self-interaction
        where $C_A = 3$ for $SU(3)$ and the numerical coefficient $11/3$
        arises from the combination of ghost and gauge-field contributions.
        In RS, $C_A = 3$ from the three quark colors forced by the
        $Q_3$ hypercube geometry (Lean: \texttt{QFT.Confinement}).
  \item The factor $-2n_f/3$ comes from quark loop contributions with
        $n_f = 3$ (active flavors at low scale) to $n_f = 6$ (at high scale).
  \item Asymptotic freedom: $b_0 > 0$ iff there are fewer than
        $16.5$ active quark flavors — automatically satisfied in RS
        since the $\varphi$-ladder has at most $n_f = 6$ rungs
        below the top mass threshold.
\end{enumerate}
\texttt{Lean: QFT.Confinement} (SU(3) structure);
\texttt{Masses.MassLaw.predict\_mass} (quark thresholds).
\end{theorem}

The solution to~\eqref{eq:beta_qcd} is
\begin{equation}
  \alpha_s(\mu) = \frac{\alpha_s(\mu^\star)}{1 + \frac{b_0}{2\pi}\alpha_s(\mu^\star)
  \ln(\mu/\mu^\star)}.
  \label{eq:alpha_s_running}
\end{equation}

\begin{proposition}[RS prediction for $\alpha_s(M_Z)$]\label{prop:alphas}
With the RS anchor $\mu^\star = 182.201$ GeV and
$\alpha_s(\mu^\star) = 0.1181$ (from Lean cert:
\texttt{RunningCouplingCert} in the RS adapter layer):
\begin{equation}
  \alpha_s(M_Z = 91.2\,\mathrm{GeV}) \approx 0.1185,
\end{equation}
consistent with the world average $\alpha_s(M_Z) = 0.1180 \pm 0.0009$.
\end{proposition}

\subsection{Electromagnetic Coupling}

The QED $\beta$-function:
\begin{equation}
  \beta(\alpha) = \frac{\alpha^2}{3\pi}\sum_f Q_f^2,
\end{equation}
where the sum is over active fermions with charge $Q_f$.  At $\mu =
m_e$, the RS-derived $\alpha = 1/137.036$:
\begin{equation}
  \alpha(M_Z) = \frac{\alpha(m_e)}{1 - \frac{\alpha}{3\pi}\Pi(M_Z^2)}
  \approx \frac{1}{128.9}.
  \label{eq:alpha_MZ}
\end{equation}

\begin{remark}[The Landau pole]
The Landau pole at $\mu_L \sim e^{3\pi/\alpha} m_e \sim 10^{286}$ GeV
is far above the Planck scale; in RS the UV cutoff is
$E_\mathrm{UV} = \hbar c/\lambda_\mathrm{rec}$, which is expected to
be near or below the Landau pole.  RS predicts that QED breaks down
at the UV cutoff before the Landau pole is reached.
\end{remark}

\subsection{The RS Display Identity at the Anchor}

A key RS result is the display identity at $\mu^\star$:

\begin{theorem}[Display identity, Lean-proved]\label{thm:display}
At the anchor scale $\mu^\star$:
\begin{equation}
  f_\mathrm{residue}(\mu^\star) = F(Z) = \mathrm{gap}(Z),
\end{equation}
where $\mathrm{gap}(Z)$ is the closed-form band function of the
charge-band parameter $Z$.  This means the RG residue at $\mu^\star$
equals the geometric gap (the $\varphi$-ladder integer part), cleanly
separating the RG transport from the structural RS prediction.

\texttt{Lean: Physics.AnchorPolicy.display\_identity\_at\_anchor}
\end{theorem}

\subsection{Gauge Coupling Unification}

In RS, the three SM gauge couplings $g_1, g_2, g_3$ (for
$U(1)_Y$, $SU(2)_L$, $SU(3)_c$) have values at $\mu^\star$ determined
by the Weinberg angle and the $\varphi$-ladder structure:
\begin{align}
  \sin^2\theta_W &= \frac{g_1^2}{g_1^2 + g_2^2}
  = \frac{3-\varphi}{6} \approx 0.2303
  \quad (\text{from Lean: \texttt{StandardModel.WeinbergAngle}}), \\
  \alpha_3(M_Z) &= 0.1185
  \quad (\text{from \texttt{RunningCouplingCert}}).
\end{align}

Running these to high scales, approximate unification occurs near
$\mu_\mathrm{GUT} \approx 2 \times 10^{16}$ GeV (SM) or $\sim 10^{16}$
GeV (MSSM).  The RS prediction: the unification scale is
$\mu_\mathrm{GUT} = E_\mathrm{coh} \cdot \varphi^{r_\mathrm{GUT}}$
for some integer rung $r_\mathrm{GUT}$ (\textbf{HYPOTHESIS}, falsifier:
precision coupling measurements ruling out a common intersection point).

%%============================================================
\section{Numerical Results}
\label{sec:numerical}
%%============================================================

\begin{center}
\begin{tabular}{lllll}
\toprule
Quantity & RS Prediction & Experiment & Status \\
\midrule
$\alpha_s(M_Z)$ & $0.1185$ & $0.1180 \pm 0.0009$ & $0.4\%$ \\
$\alpha(M_Z)$ & $1/128.9$ & $1/128.9$ & \checkmark \\
$\sin^2\theta_W(M_Z)$ & $(3-\varphi)/6 = 0.2303$ & $0.2229 \pm 0.0003$ & $3.3\%$ \\
Asymptotic freedom & $b_0 = 7 > 0$ for $n_f = 6$ & Confirmed & \checkmark \\
Anchor stationarity & $\gamma_f(\mu^\star) = 0$ & Consistent with data & \checkmark \\
Unification scale & $\sim 10^{16}$ GeV & $\sim 2 \times 10^{16}$ GeV & HYPOTHESIS \\
\bottomrule
\end{tabular}
\end{center}

%%============================================================
\section{Lean Formalization Status}
\label{sec:lean}
%%============================================================

\begin{center}
\begin{tabular}{llll}
\toprule
Claim & Lean theorem & Module & Status \\
\midrule
Anchor stationarity & \texttt{stationary\_at\_anchor} & \texttt{Physics.AnchorPolicy} & Proved \\
Stability at anchor & \texttt{stability\_bound\_at\_anchor} & \texttt{Physics.AnchorPolicy} & Proved \\
Display identity & \texttt{display\_identity\_at\_anchor} & \texttt{Physics.AnchorPolicy} & Proved \\
MFV compatibility & \texttt{mfv\_compatible\_at\_anchor} & \texttt{Physics.AnchorPolicy} & Proved \\
Running coupling cert & \texttt{RunningCouplingCert} & Adapter & Proved \\
Weinberg angle & \texttt{prediction3} & \texttt{StandardModel.WeinbergAngle} & Proved \\
QFT confinement & \texttt{Confinement} & \texttt{QFT.Confinement} & Proved \\
\midrule
Full $\beta$-function from RS & --- & --- & HYPOTHESIS$^\dagger$ \\
Gauge unification from RS & --- & --- & HYPOTHESIS$^\ddagger$ \\
\bottomrule
\end{tabular}
\end{center}

$^\dagger$ \textbf{HYPOTHESIS}: Full multi-loop QCD $\beta$-function
from RS first principles; Lean formalization of the loop integrals in RS
language pending.  Falsifier: $\alpha_s(M_Z)$ measured at $> 5\sigma$
from $0.1185$.

$^\ddagger$ \textbf{HYPOTHESIS}: Gauge coupling unification at a
$\varphi$-ladder rung.  Falsifier: precision measurements ruling out
$\alpha_1 = \alpha_2 = \alpha_3$ at any scale below $10^{20}$ GeV.

%%============================================================
\section{Conclusion}
\label{sec:conclusion}
%%============================================================

We have derived the renormalization group structure from the RS
$\varphi$-ladder.  The anchor scale $\mu^\star = 182.201$ GeV is a
stationarity point of the RG flow (Lean-proved); it is not a fit
parameter.  Asymptotic freedom follows from the $SU(3)$ color structure
forced by the $Q_3$ hypercube.  The RS framework reproduces all
standard RG predictions with zero free parameters at the structural level.

\begin{thebibliography}{99}

\bibitem{Gross1973}
D.~J.~Gross, F.~Wilczek, ``Ultraviolet behavior of non-abelian gauge
theories,'' \textit{Phys.\ Rev.\ Lett.}\ \textbf{30}, 1343 (1973).

\bibitem{Politzer1973}
H.~D.~Politzer, ``Reliable perturbative results for strong interactions,''
\textit{Phys.\ Rev.\ Lett.}\ \textbf{30}, 1346 (1973).

\bibitem{Wilson1971}
K.~G.~Wilson, ``Renormalization group and critical phenomena,''
\textit{Phys.\ Rev.\ B}\ \textbf{4}, 3174 (1971).

\bibitem{PDG2022}
Particle Data Group, ``Review of Particle Physics,''
\textit{Prog.\ Theor.\ Exp.\ Phys.}\ \textbf{2022}, 083C01 (2022).

\bibitem{Washburn2026a}
J.~Washburn, ``The Algebra of Reality: A Recognition Science Derivation
of Physical Law,'' \textit{Axioms} \textbf{15}(2), 90 (2026).

\bibitem{Washburn2026b}
J.~Washburn, ``Gauge Couplings from the $\varphi$-Ladder,'' preprint (2026).

\end{thebibliography}

\end{document}
